% GENERATED FILE. Edit canonical lessons, then run npm run generate:books.
% canonical-source-hash: fnv1a64:d769c8a1f2b99e8c
% canonical-lessons: PA-C357-rabar, PA-C357-kapi, PA-C357-siahi, PA-C357-salet, PA-C357-panna, PA-R357-review-351-357

\chapter{My School Things}
\label{ch:pa-a2-my-school-things}
\hlchaptermodality{357}

\begin{chapteropening}
\textbf{By the end of this chapter:} \emph{I can say what I have with me for school.}

Complete the last lesson of chapter 357: I can say what I have with me for school.
\end{chapteropening}

\section[\texorpdfstring{\pa{ਰਬੜ} (rabaṛ)}{ਰਬੜ (rabaṛ)}]{\pa{ਰਬੜ} (rabaṛ) — an eraser; rubber}

\label{lesson:PA-C357-rabar}



\begin{quote}
Before the new one: say the Punjabi for a photo album, then the Punjabi for a picture frame.
\end{quote}

\subsection*{You'll want to know: \pa{ਰਬੜ}}
\textbf{\pa{ਰਬੜ}} — \emph{rabaṛ} — \textquotedblleft{}an eraser; rubber\textquotedblright{}.

\begin{quote}
\pa{ਮੇਰੇ} \pa{ਕੋਲ} \pa{ਇੱਕ} \pa{ਨਵਾਂ} \pa{ਰਬੜ} \pa{ਹੈ।}
\end{quote}

\begin{quote}
\emph{mere kol ikk navāṁ rabaṛ hai.}
\end{quote}

\begin{quote}
I have a new eraser.
\end{quote}

A masculine noun.

\begin{cousinweb}
\textbf{\pa{ਰਬੜ}} (\emph{rabaṛ}) is borrowed from English \emph{rubber}.
\end{cousinweb}

\begin{grammarlens}[title={What you've built}]
One more word for this chapter.
\end{grammarlens}

\subsection*{Your turn}
\begin{itemize}
\raggedright
  \item \emph{Say it:} \emph{rabaṛ}
  \item \emph{Say it:} \emph{rabaṛ}, once more
  \item \emph{Say it:} \emph{mere kol ikk navāṁ rabaṛ hai.} — I have a new eraser.
  \item \emph{Recall:} say the Punjabi for a photo album, then the Punjabi for a picture frame, then say \emph{rabaṛ} again
\end{itemize}

\begin{tcolorbox}[breakable,colback=teal!4,colframe=teal!35!black,title={Before you move on}]
What does \pa{ਰਬੜ} mean? (\textquotedblleft{}An eraser; rubber\textquotedblright{}.) Say it once more.
\end{tcolorbox}
\section[\texorpdfstring{\pa{ਕਾਪੀ} (kāpī)}{ਕਾਪੀ (kāpī)}]{\pa{ਕਾਪੀ} (kāpī) — a notebook; an exercise book}

\label{lesson:PA-C357-kapi}



\begin{quote}
Before the new one: say the Punjabi for a picture frame, then the Punjabi for an eraser.
\end{quote}

\subsection*{You'll want to know: \pa{ਕਾਪੀ}}
\textbf{\pa{ਕਾਪੀ}} — \emph{kāpī} — \textquotedblleft{}a notebook; an exercise book\textquotedblright{}.

\begin{quote}
\pa{ਮੇਰੇ} \pa{ਬਸਤੇ} \pa{ਵਿੱਚ} \pa{ਤਿੰਨ} \pa{ਕਾਪੀਆਂ} \pa{ਹਨ।}
\end{quote}

\begin{quote}
\emph{mere baste vich tinn kāpīāṁ han.}
\end{quote}

\begin{quote}
There are three exercise books in my school bag.
\end{quote}

A feminine noun.

\begin{cousinweb}
\textbf{\pa{ਕਾਪੀ}} (\emph{kāpī}) is borrowed from English \emph{copy}.
\end{cousinweb}

\begin{grammarlens}[title={What you've built}]
One more word for this chapter.
\end{grammarlens}

\subsection*{Your turn}
\begin{itemize}
\raggedright
  \item \emph{Say it:} \emph{kāpī}
  \item \emph{Say it:} \emph{kāpī}, once more
  \item \emph{Say it:} \emph{mere baste vich tinn kāpīāṁ han.} — There are three exercise books in my school bag.
  \item \emph{Recall:} say the Punjabi for a picture frame, then the Punjabi for an eraser, then say \emph{kāpī} again
  \item \emph{Write it:} copy \textbf{\pa{ਕਾਪੀ}} once, checking the whole word
\end{itemize}

\begin{tcolorbox}[breakable,colback=teal!4,colframe=teal!35!black,title={Before you move on}]
What does \pa{ਕਾਪੀ} mean? (\textquotedblleft{}A notebook; an exercise book\textquotedblright{}.) Say it once more.
\end{tcolorbox}
\section[\texorpdfstring{\pa{ਸਿਆਹੀ} (siāhī)}{ਸਿਆਹੀ (siāhī)}]{\pa{ਸਿਆਹੀ} (siāhī) — ink}

\label{lesson:PA-C357-siahi}



\begin{quote}
Before the new one: say the Punjabi for an eraser, then the Punjabi for a notebook.
\end{quote}

\subsection*{You'll want to know: \pa{ਸਿਆਹੀ}}
\textbf{\pa{ਸਿਆਹੀ}} — \emph{siāhī} — \textquotedblleft{}ink\textquotedblright{}.

\begin{quote}
\pa{ਮੇਰੀ} \pa{ਕਲਮ} \pa{ਵਿੱਚ} \pa{ਨੀਲੀ} \pa{ਸਿਆਹੀ} \pa{ਹੈ।}
\end{quote}

\begin{quote}
\emph{merī kalam vich nīlī siāhī hai.}
\end{quote}

\begin{quote}
There is blue ink in my pen.
\end{quote}

A feminine noun.

\begin{cousinweb}
\textbf{\pa{ਸਿਆਹੀ}} (\emph{siāhī}) is a loan, usually traced to Persian \emph{siyāhī}, \textquotedblleft{}blackness\textquotedblright{}.
\end{cousinweb}

\begin{grammarlens}[title={What you've built}]
One more word for this chapter.
\end{grammarlens}

\subsection*{Your turn}
\begin{itemize}
\raggedright
  \item \emph{Say it:} \emph{siāhī}
  \item \emph{Say it:} \emph{siāhī}, once more
  \item \emph{Say it:} \emph{merī kalam vich nīlī siāhī hai.} — There is blue ink in my pen.
  \item \emph{Recall:} say the Punjabi for an eraser, then the Punjabi for a notebook, then say \emph{siāhī} again
  \item \emph{Write it:} copy \textbf{\pa{ਸਿਆਹੀ}} once, checking the whole word
\end{itemize}

\begin{tcolorbox}[breakable,colback=teal!4,colframe=teal!35!black,title={Before you move on}]
What does \pa{ਸਿਆਹੀ} mean? (\textquotedblleft{}Ink\textquotedblright{}.) Say it once more.
\end{tcolorbox}
\section[\texorpdfstring{\pa{ਸਲੇਟ} (saleṭ)}{ਸਲੇਟ (saleṭ)}]{\pa{ਸਲੇਟ} (saleṭ) — a slate (to write on)}

\label{lesson:PA-C357-salet}



\begin{quote}
Before the new one: say the Punjabi for a notebook, then the Punjabi for ink.
\end{quote}

\subsection*{You'll want to know: \pa{ਸਲੇਟ}}
\textbf{\pa{ਸਲੇਟ}} — \emph{saleṭ} — \textquotedblleft{}a slate (to write on)\textquotedblright{}.

\begin{quote}
\pa{ਛੋਟੇ} \pa{ਬੱਚੇ} \pa{ਸਲੇਟ} \pa{ਉੱਤੇ} \pa{ਲਿਖਦੇ} \pa{ਹਨ।}
\end{quote}

\begin{quote}
\emph{chhoṭe bachche saleṭ utte likhde han.}
\end{quote}

\begin{quote}
Little children write on a slate.
\end{quote}

A feminine noun.

\begin{cousinweb}
\textbf{\pa{ਸਲੇਟ}} (\emph{saleṭ}) is borrowed from English \emph{slate}.
\end{cousinweb}

\begin{grammarlens}[title={What you've built}]
One more word for this chapter.
\end{grammarlens}

\subsection*{Your turn}
\begin{itemize}
\raggedright
  \item \emph{Say it:} \emph{saleṭ}
  \item \emph{Say it:} \emph{saleṭ}, once more
  \item \emph{Say it:} \emph{chhoṭe bachche saleṭ utte likhde han.} — Little children write on a slate.
  \item \emph{Recall:} say the Punjabi for a notebook, then the Punjabi for ink, then say \emph{saleṭ} again
  \item \emph{Write it:} copy \textbf{\pa{ਸਲੇਟ}} once, checking the whole word
\end{itemize}

\begin{tcolorbox}[breakable,colback=teal!4,colframe=teal!35!black,title={Before you move on}]
What does \pa{ਸਲੇਟ} mean? (\textquotedblleft{}A slate (to write on)\textquotedblright{}.) Say it once more.
\end{tcolorbox}
\section[\texorpdfstring{\pa{ਪੰਨਾ} (pannā)}{ਪੰਨਾ (pannā)}]{\pa{ਪੰਨਾ} (pannā) — a page; a leaf of a book}

\label{lesson:PA-C357-panna}



\begin{quote}
Before the new one: say the Punjabi for ink, then the Punjabi for a slate.
\end{quote}

\subsection*{You'll want to know: \pa{ਪੰਨਾ}}
\textbf{\pa{ਪੰਨਾ}} — \emph{pannā} — \textquotedblleft{}a page; a leaf of a book\textquotedblright{}.

\begin{quote}
\pa{ਮੇਰੀ} \pa{ਕਿਤਾਬ} \pa{ਵਿੱਚ} \pa{ਸੌ} \pa{ਪੰਨੇ} \pa{ਹਨ।}
\end{quote}

\begin{quote}
\emph{merī kitāb vich sau panne han.}
\end{quote}

\begin{quote}
My book has a hundred pages.
\end{quote}

A masculine noun.

\begin{cousinweb}
\textbf{\pa{ਪੰਨਾ}} (\emph{pannā}) is inherited from Sanskrit \emph{parṇa}, \textquotedblleft{}plumage, foliage\textquotedblright{}; CDIAL lists Punjabi \emph{pannā}, \textquotedblleft{}leaf, page\textquotedblright{}, under it (CDIAL 7918).
\end{cousinweb}

\begin{grammarlens}[title={What you've built}]
One more word for this chapter.
\end{grammarlens}

\subsection*{Your turn}
\begin{itemize}
\raggedright
  \item \emph{Say it:} \emph{pannā}
  \item \emph{Say it:} \emph{pannā}, once more
  \item \emph{Say it:} \emph{merī kitāb vich sau panne han.} — My book has a hundred pages.
  \item \emph{Recall:} say the Punjabi for ink, then the Punjabi for a slate, then say \emph{pannā} again
  \item \emph{Write it:} copy \textbf{\pa{ਪੰਨਾ}} once, checking the whole word
\end{itemize}

\begin{tcolorbox}[breakable,colback=teal!4,colframe=teal!35!black,title={Before you move on}]
What does \pa{ਪੰਨਾ} mean? (\textquotedblleft{}A page; a leaf of a book\textquotedblright{}.) Say it once more.
\end{tcolorbox}
\section[(dialogue)]{Looking back: from Ball Games to My School Things}

\label{lesson:PA-R357-review-351-357}



\begin{quote}
Say the Punjabi for a ball and hockey.
\end{quote}

\subsection*{Your turn}
Say these aloud:

\begin{itemize}
\raggedright
  \item a ball, hockey, a team, a match, to defeat
  \item chess, a round of a game, one's turn, a game piece, a prize
  \item a spinning top, a contest, to leap, to make something fly, to play with a child
  \item a story, to tell a story, a tale, a queen, magic
  \item a riddle, to solve a riddle, a joke, to make someone laugh, to tease
  \item a camera, a picture, a photograph, a photo album, a picture frame
  \item an eraser, a notebook, ink, a slate, a page
\end{itemize}

\begin{tcolorbox}[breakable,colback=teal!4,colframe=teal!35!black,title={Before you move on}]
A ball? (\textbf{\pa{ਗੇਂਦ}}, \emph{gend}.) Hockey? (\textbf{\pa{ਹਾਕੀ}}, \emph{hākī}.) A team? (\textbf{\pa{ਟੀਮ}}, \emph{ṭīm}.) A match? (\textbf{\pa{ਮੈਚ}}, \emph{maich}.) To defeat? (\textbf{\pa{ਹਰਾਉਣਾ}}, \emph{harāuṇā}.) Chess? (\textbf{\pa{ਸ਼ਤਰੰਜ}}, \emph{shatranj}.) A round of a game? (\textbf{\pa{ਬਾਜ਼ੀ}}, \emph{bāzī}.) One's turn? (\textbf{\pa{ਵਾਰੀ}}, \emph{vārī}.) A game piece? (\textbf{\pa{ਗੋਟੀ}}, \emph{goṭī}.) A prize? (\textbf{\pa{ਇਨਾਮ}}, \emph{inām}.) A spinning top? (\textbf{\pa{ਲਾਟੂ}}, \emph{lāṭū}.) A contest? (\textbf{\pa{ਮੁਕਾਬਲਾ}}, \emph{mukāblā}.) To leap? (\textbf{\pa{ਕੁੱਦਣਾ}}, \emph{kuddṇā}.) To make something fly? (\textbf{\pa{ਉਡਾਉਣਾ}}, \emph{uḍāuṇā}.) To play with a child? (\textbf{\pa{ਖਿਡਾਉਣਾ}}, \emph{khiḍāuṇā}.) A story? (\textbf{\pa{ਕਹਾਣੀ}}, \emph{kahāṇī}.) To tell a story? (\textbf{\pa{ਸੁਣਾਉਣਾ}}, \emph{suṇāuṇā}.) A tale? (\textbf{\pa{ਕਿੱਸਾ}}, \emph{kissā}.) A queen? (\textbf{\pa{ਰਾਣੀ}}, \emph{rāṇī}.) Magic? (\textbf{\pa{ਜਾਦੂ}}, \emph{jādū}.) A riddle? (\textbf{\pa{ਪਹੇਲੀ}}, \emph{pahelī}.) To solve a riddle? (\textbf{\pa{ਬੁੱਝਣਾ}}, \emph{bujjhṇā}.) A joke? (\textbf{\pa{ਮਜ਼ਾਕ}}, \emph{mazāk}.) To make someone laugh? (\textbf{\pa{ਹਸਾਉਣਾ}}, \emph{hasāuṇā}.) To tease? (\textbf{\pa{ਛੇੜਨਾ}}, \emph{chheṛnā}.) A camera? (\textbf{\pa{ਕੈਮਰਾ}}, \emph{kaimrā}.) A picture? (\textbf{\pa{ਤਸਵੀਰ}}, \emph{tasvīr}.) A photograph? (\textbf{\pa{ਫ਼ੋਟੋ}}, \emph{foṭo}.) A photo album? (\textbf{\pa{ਐਲਬਮ}}, \emph{ailbam}.) A picture frame? (\textbf{\pa{ਫ਼ਰੇਮ}}, \emph{farem}.) An eraser? (\textbf{\pa{ਰਬੜ}}, \emph{rabaṛ}.) A notebook? (\textbf{\pa{ਕਾਪੀ}}, \emph{kāpī}.) Ink? (\textbf{\pa{ਸਿਆਹੀ}}, \emph{siāhī}.) A slate? (\textbf{\pa{ਸਲੇਟ}}, \emph{saleṭ}.) A page? (\textbf{\pa{ਪੰਨਾ}}, \emph{pannā}.)
\end{tcolorbox}
